\begin{longtable}{@{}p{1.1cm}p{4.6cm}ccccc>{\raggedleft\arraybackslash}p{1.5cm}@{}}
\caption{Penilaian alternatif pendekatan menggunakan \emph{Weighted Sum Model}.}
\label{tab:wsm} \\
\toprule
\textbf{Kode} & \textbf{Alternatif pendekatan} & \textbf{C1} & \textbf{C2} & \textbf{C3} & \textbf{C4} & \textbf{C5} & \textbf{Skor} \\
& \emph{Bobot} & 0,25 & 0,25 & 0,20 & 0,20 & 0,10 & \\
\midrule
\endfirsthead

\multicolumn{8}{@{}l}{\small\emph{Tabel \thetable{} (lanjutan)}} \\
\toprule
\textbf{Kode} & \textbf{Alternatif pendekatan} & \textbf{C1} & \textbf{C2} & \textbf{C3} & \textbf{C4} & \textbf{C5} & \textbf{Skor} \\
\midrule
\endhead

\bottomrule
\endlastfoot

P1 & Pemantauan manual berbasis ambang & 1 & 1 & 3 & 5 & 3 & 2,40 \\[2pt]

P2 & Model prediksi tanpa penelusuran & 5 & 1 & 1 & 4 & 2 & 2,70 \\[2pt]

P3 & Model bahasa besar tanpa penelusuran & 1 & 2 & 2 & 3 & 1 & 1,85 \\[2pt]

P4 & RAG dengan kueri statis & 1 & 2 & 5 & 3 & 4 & 2,75 \\[2pt]

P5 & Model fisiologis mekanistik & 5 & 3 & 3 & 1 & 3 & 3,10 \\[2pt]

P6 & \textbf{RAG dengan kueri terkondisi-prediksi} & 5 & 5 & 5 & 3 & 4 & \textbf{4,50} \\

\end{longtable}

\noindent\tabnotestyle Skor diberikan pada skala 1 sampai 5, dengan 1 menandakan kesesuaian sangat rendah dan 5 kesesuaian sangat tinggi. Kolom skor memuat hasil penjumlahan berbobot menurut Persamaan~\ref{eq:wsm}. Seluruh kriteria dinilai pada skala yang sama, sehingga syarat kesebandingan satuan yang dinyatakan \textcite{triantaphyllou1989} terpenuhi.\normalsize
